\section{Introduction}
\label{sec:intro}

\Cref{fig:example} shows the failure this report studies. The model is given a two-line memory, asked where Ben
lives, and told to write a short note before its answer. The note quotes the relevant memory line, ``Ben lives in
Manchester'', and the answer line then says ``unknown''.

\begin{figure}[t]
\centering\small
\fbox{\begin{minipage}{0.95\linewidth}
\begin{tabular}{@{}l@{\hspace{0.8em}}>{\raggedright\arraybackslash}p{0.76\linewidth}@{}}
\textbf{Memory} & \texttt{user, 2023/03/01 (Wed) 10:00: Ben lives in Manchester.}\\
 & \texttt{user, 2023/03/02 (Thu) 10:00: Ben works as a manager.}\\[3pt]
\textbf{Question} & \texttt{(asked on 2023/03/03 (Fri) 12:00) Where does Ben live?}\\[3pt]
\textbf{Unedited model} & \texttt{Note: user mentioned that Ben lives in Manchester on 2023/03/01 (Wed) 10:00.}\\
 & \texttt{Answer: unknown}\\[3pt]
\textbf{Edited model} & \texttt{Note: user mentioned that Ben lives in Manchester on 2023/03/01 (Wed) 10:00.}\\
 & \texttt{Answer: Manchester}
\end{tabular}
\end{minipage}}
\caption{A false abstention by Qwen2.5-3B-Instruct, copied verbatim from stored outputs (case
\texttt{E-DE01\_ANSa\_H0}). The prompt, given in full in \cref{app:prompt}, ends with ``Note:'', so the
model's output begins with the note. The unedited model quotes the answer
in its note and then abstains. The edited output is the agent's layer-35 edit at dose~1 with the unedited model's
note held fixed (the controlled-prefix setting of \cref{sec:evaluation}). The example comes from the exploration
split, which the agent studied, so it illustrates the failure and is not evidence for the repair.}
\label{fig:example}
\end{figure}

We call this a \emph{false abstention}: the context supports an answer, but the model declines to give one. An
\emph{appropriate abstention} is different, because there the context really does not contain the answer. In this
case study the pattern is common. On the sealed confirmation split, the unedited model abstained falsely on 261
answerable questions, and on 232 of them its note named the correct city, by a frozen matching rule, before the final ``unknown''. The note is
a prompted summary, not a faithful record of the model's computation, so it does not show that the model ``knew''
the answer (\cref{sec:setup}). It does show that the answer was available in the model's own output a few tokens
before the model declined to give it.

\paragraph{Why the failure is hard to fix.}
Recovering these answers is easy if one is willing to make the model answer everything. A useful fix has to release
the missed answers while keeping the model's abstentions on questions that its context cannot answer. Earlier
experiments on the same model and task family (summarised in \cref{sec:background}) found that simple interventions
do not achieve this. Adding a content-free activation direction at layer 29 recovered answers but raised the
unsupported-answer rate on unanswerable questions to 0.188. Removing the ``unknown'' option from the prompt raised it
to 0.451. We therefore score every intervention on both sides of the trade-off: the share of false abstentions it
recovers, and the share of unanswerable questions it answers anyway.

\paragraph{Why evaluate agent-proposed weight edits.}
Recent work shows that a language-model agent can inspect another network and propose a small weight edit that
changes its behaviour \citep{takehi2026wsp}. A weight edit is persistent and cheap to apply. Unlike a prompt change,
it can be inspected and compared directly with other interventions. When the agent is required to state a hypothesis
before it edits, that hypothesis also makes predictions that can be tested. Two risks come with this. An agent that
studies a split of the data can fit that split, and its stated explanation can be wrong even when its edit works.
Evaluating such an edit therefore needs data the agent never saw, a fixed analysis plan, comparison with simpler
interventions, and separate tests for the edit and for the explanation.

\paragraph{What we did.}
A Claude Opus~5.5 agent inspected Qwen2.5-3B-Instruct on an exploration split and submitted three weight edits, each
with a written hypothesis and predictions (\cref{sec:proposals}). We then evaluated the edits in four stages
(\cref{fig:pipeline}). The exploration data screened the submissions; a separate selection split chose a repair
candidate and its dose; the candidates, doses, comparator settings and analysis code were frozen by commit; and only
then were the sealed confirmation and transfer splits generated and read, once each. The evaluation asks two separate
questions. The first is whether an edit repairs false abstention better than five frozen comparators: logit bias, a
prompt instruction, option removal, activation steering and directional orthogonalisation. The second is whether the
agent's proposed pathway and explanation survive intervention tests.

\paragraph{What we found.}
The selected repair is a rank-1 update to the attention output projection of layer 35, the model's last layer. On
sealed data it met the pre-specified repair criterion and had a better recovery-versus-unsupported-answer trade-off
than the prompt instruction, activation steering and orthogonalisation (\cref{sec:results}). It did not clearly
beat the simplest baseline, a logit bias against ``unknown''. That comparison was inconclusive, and removing the
``unknown'' option recovered more answers at a far higher rate of unsupported answers. The edit also broke 3\% of
initially correct answers, which logit bias and steering did not. A different submission, a two-head edit, showed that
changing one layer-33 component is sufficient to release answers. The frozen tests did not establish the agent's
explanation of what that component computes (\cref{sec:explain}).

\paragraph{Contribution.}
The study is a controlled case study with one model, one task family and one agent run. Its contribution is the
evaluation, not the finding that agents can edit weights. We make three points:
\begin{enumerate}[leftmargin=*]
  \item We give an evaluation design for agent-proposed edits that keeps the agent's data, the selection data and
        the confirmation data apart, freezes the analysis before confirmation, and tests repair and explanation
        separately (\cref{sec:evaluation}).
  \item Under that design, an agent-proposed edit repaired false abstention with the evidence in view, and improved
        on three of five simpler interventions under corrected tests. The comparison with logit bias, the simplest
        output-level fix, remains open (\cref{sec:results}).
  \item The edit that worked and the explanation offered for it were separable. A pathway test passed while the
        agent's account of the same component was not established (\cref{sec:explain}).
\end{enumerate}
\Cref{sec:related} places these results in the literature. \Cref{sec:discussion,sec:limits} discuss what they do
not show.
